\section{Experiments}
\label{sec:experiments}

We test whether controllers synthesized from calibration-time dynamics remain
reliable under prompt, adversarial, language, and context-length shifts.  All
identified dynamics, disturbance channels, gains, and hyperparameters are
frozen before shifted evaluation.  A-LQR and $\mathrm{H}_{\infty}$ use the
same state coordinates, references, intervention channels, and quadratic
weights, isolating the effect of nominal versus worst-case synthesis.  Full
protocols, per-condition results, and standard errors are reported in
Appendix~\ref{app:full-results}.

\begin{figure*}[t]
    \centering
    \includegraphics[width=\textwidth]
    {Figures/benchmark-evaluation-overview.pdf}
    \caption{
        \textbf{Evaluation under distribution shift.}
        Controllers are calibrated in-distribution and frozen before testing.
        \emph{Aligned} benchmarks directly score the steered behavior;
        \emph{orthogonal} benchmarks require attaining a separate steering
        target while preserving task capability.
    }
    \label{fig:benchmark-overview}
\end{figure*}

Figure~\ref{fig:benchmark-overview} summarizes the four evaluation settings.
We use ``orthogonal'' to describe the roles of the objectives, not geometric
orthogonality between their representation directions.


% ============================================================
% TruthfulQA
% ============================================================

\subsection{Prompt Shift: TruthfulQA}
\label{sec:exp-truthfulqa}

We calibrate on English prompts and test matched English inputs (ID) and
Spanish-input, English-output transfer (OOD).  We measure the joint
truthfulness--quality trade-off by
$\mathrm{QAT}=\mathrm{True}(\%)\min\{\mathrm{IR}/2,\mathrm{F}/2\}$, where
IR and F denote instruction relevance and fluency.

\begin{figure*}[t]
    \centering
    \begin{subfigure}[t]{0.49\textwidth}
        \centering
        \includegraphics[width=\linewidth]
        {Figures/truthfulqa-id-quality-frontier.pdf}
        \caption{Matched English evaluation.}
        \label{fig:truthfulqa-id}
    \end{subfigure}
    \hfill
    \begin{subfigure}[t]{0.49\textwidth}
        \centering
        \includegraphics[width=\linewidth]
        {Figures/truthfulqa-ood-quality-frontier.pdf}
        \caption{Spanish-to-English transfer.}
        \label{fig:truthfulqa-ood}
    \end{subfigure}
    \caption{
        \textbf{Truthfulness--quality trade-off.}
        Colors denote models, markers denote methods, and dashed contours show
        equal QAT; diamonds mark $\mathrm{H}_{\infty}$.
    }
    \label{fig:truthfulqa-quality-tradeoff}
\end{figure*}

Under OOD transfer, $\mathrm{H}_{\infty}$ attains the best QAT on both models
(Figure~\ref{fig:truthfulqa-quality-tradeoff}): $39.50$ on Gemma-2-2B versus
$39.12$ for A-LQR, and $41.33$ on Llama-3-8B versus $35.29$.  The ID Llama
result provides an important qualification: robust synthesis substantially
raises truthfulness but not QAT because the stronger intervention reduces
response quality.  Its advantage is therefore a better frontier under shift,
not uniform dominance of every nominal metric.


% ============================================================
% HarmBench
% ============================================================

\subsection{Adversarial Prompt Shift: HarmBench}
\label{sec:exp-harmbench}

We evaluate refusal on direct HarmBench requests and five frozen jailbreak
templates over the same 240 behaviors.  ASR is the primary endpoint, while
safe-concept relevance checks that lower ASR reflects the intended refusal.

\begin{figure}[t]
    \centering
    \includegraphics[width=\linewidth]
    {Figures/harmbench-robust-refusal.pdf}
    \caption{
        \textbf{Robust refusal under adversarial prompts.}
        \textbf{(A)} ASR across direct and jailbreak conditions.
        \textbf{(B)} Mean ASR versus safe-concept relevance.
    }
    \label{fig:harmbench-refusal}
\end{figure}

Across all six conditions, $\mathrm{H}_{\infty}$ keeps ASR below $3\%$
(Figure~\ref{fig:harmbench-refusal}).  Mean ASR falls from $6.18\%$ with
A-LQR to $1.25\%$---an $80\%$ reduction---while safe-concept relevance
slightly increases from $1.968$ to $1.977$.  The gain therefore reflects more
reliable refusal rather than abandonment of the safety target.


% ============================================================
% MGSM
% ============================================================

\subsection{Language Shift: MGSM}
\label{sec:exp-mgsm}

We calibrate in English and transfer the frozen controllers to five unseen
MGSM languages.  Joint target attainment and capability retention are measured
by $Q_{\mathrm{tr}}=\min\{T/2,\mathrm{IR}/2,\mathrm{F}/2\}$ and
$\mathrm{QAA}=\mathrm{Acc}(\%)Q_{\mathrm{tr}}$, where $T$ is target-language
relevance.

\begin{figure}[t]
    \centering
    \includegraphics[width=\linewidth]
    {Figures/mgsm-transfer-overview.pdf}
    \caption{
        \textbf{Cross-lingual transfer on MGSM.}
        Left: capability--steering frontier with equal-QAA contours.
        Right: accuracy by unseen language.  Calibration uses English only.
    }
    \label{fig:mgsm-transfer}
\end{figure}

$\mathrm{H}_{\infty}$ leads the steering baselines in four of five languages
(Figure~\ref{fig:mgsm-transfer}), retaining $32.2\%$ mean accuracy versus
$22.4\%$ for A-LQR and $0\%$ for S-PID.  This is not obtained by weakening
the intervention: $Q_{\mathrm{tr}}$ rises from $0.542$ to $0.625$, increasing
QAA from $12.14$ to $20.13$.  Although the unmodified model has higher raw
accuracy, its zero target-language relevance means that it does not achieve
the steering objective.


% ============================================================
% Context length
% ============================================================

\subsection{Context-Length Shift: L-CiteEval}
\label{sec:exp-context}

We calibrate at $8$K and evaluate the frozen controllers as context length
increases, reporting answer recall (AR), citation F1 (CF1), and overall
steering quality (OS).

\begin{figure}[t]
    \centering
    \includegraphics[width=\linewidth]
    {Figures/lciteeval-context-robustness.pdf}
    \caption{
        \textbf{Robustness to context-length shift on L-CiteEval.}
        AR, CF1, and OS for controllers calibrated at $8$K and evaluated at
        increasing context lengths.
    }
    \label{fig:context-shift}
\end{figure}

The robust controller need not dominate at the matched $8$K condition, but the
ordering changes under shift (Figure~\ref{fig:context-shift}).  At $16$K,
A-LQR's CF1 falls from $5.6$ to $3.0$ and S-PID's AR from $72.1\%$ to
$61.9\%$; $\mathrm{H}_{\infty}$ instead preserves AR
($66.4\%\!\rightarrow\!65.2\%$) while CF1 rises from $6.4$ to $7.8$, yielding
the strongest controlled result.  This crossover is the expected signature of
robust synthesis: its advantage emerges as deployment dynamics depart from
calibration.


% The calibration-time predictive analysis is a central paper claim and should
% return here once its cross-model figure is complete.  Use one figure containing
% (i) S_rob versus held-out reliability and (ii) residual-subspace coverage or
% reference deviation versus certificate error.  Do not expose a visible
% ``results forthcoming'' subsection in the submitted manuscript.
